% ECONOMICS_PROBLEM: {"id": "I38-339", "title": "Cash transfers at economy-wide scale", "jel_code": "I38", "source_id": "OP-0372"}
% !TeX program = xelatex
\documentclass[11pt,a4paper]{article}
\usepackage[margin=23mm]{geometry}
\usepackage{fontspec}
\setmainfont{Latin Modern Roman}
\usepackage{amsmath,amssymb,mathtools}
\usepackage{xcolor,enumitem,booktabs,longtable,array,xurl,hyperref}
\definecolor{muted}{HTML}{53636C}
\hypersetup{colorlinks=true,urlcolor=blue,linkcolor=blue}
\setlength{\parindent}{0pt}
\setlength{\parskip}{5pt}
\newcommand{\R}{\mathbb R}
\newcommand{\E}{\mathbb E}
\newcommand{\N}{\mathbb N}
\newcommand{\ind}{\mathbf 1}
\newcommand{\Var}{\operatorname{Var}}
\newcommand{\Cov}{\operatorname{Cov}}
\newcommand{\supp}{\operatorname{supp}}
\DeclareMathOperator*{\argmin}{arg\,min}
\DeclareMathOperator*{\argmax}{arg\,max}
\newcommand{\entrymeta}[3]{{\sffamily\small\color{muted}\textbf{\texttt{#1}}\quad|\quad #2\par #3\par}}
\newcommand{\Problemref}[1]{\texttt{#1}}
\newcommand{\JELref}[2]{\texttt{#2} (JEL \texttt{#1})}
\begin{document}
\section*{Cash transfers at economy-wide scale}
\textbf{Problem ID:} \texttt{I38-339}\quad \textbf{JEL:} \texttt{I38}\par
% BEGIN ECONOMICS STATEMENT
\label{problem:OP-0372}
{\sffamily\small\color{muted}Original subject group: Development and poverty\par}
\label{definitions:problem:30007654}
\textbf{Audit classification:} Background; proposed extension. Original metadata checked and 2024 primary PDF read. NBER search metadata reports an August 2026 revision, whose full text was inaccessible; no claim is made to have inspected that revision.
\subsection*{Definitions and precise research target}
\textbf{Complete editorial problem.} Consider a specified national or regional economy with eligible and ineligible households, firms, and government. Policy \(p=(b,q,f)\) gives transfer amount \(b\) to coverage share \(q\) and uses financing rule \(f\), which explicitly specifies taxes, borrowing, aid, or spending displacement. Let \(C_{it}(p)\) be household real consumption, \(P_t(p)\) a price index, and \(Y_t(p)\) real aggregate production. Define \(\tau_C(T,p)=E[\sum_i\omega_i\{C_{iT}(p)-C_{iT}(0)\}]\) for declared population weights \(\omega_i\), with analogous price and output responses. Estimate how these effects vary across transfer intensity, supply capacity, financing, and time since disbursement. Include effects on nonrecipients, migration, firm entry, and government budgets. Randomized household transfers identify local effects only under an appropriate interference design; extrapolation to full coverage requires a validated equilibrium model or policy variation at larger scales. Determine which productivity and consumption changes survive after transfers cease. The target is a family of context-specific response functions, not one universally positive cash-transfer multiplier. Existing evidence of particular general-equilibrium effects should be treated as partial answers; an explicit source declaring this whole catalogue formulation open was not verified.

\textbf{Definitions and admissible scope.} The transfer regime fixes eligible households, amount and duration, announcement, financing, and coverage selection. Let \(q\) alter recipient saturation through a declared assignment law. The price index uses fixed baseline consumption weights; real household consumption uses comparable deflators and includes nonrecipients. Population weights sum to one. Define \(\tau_C\) for the baseline population rather than conditioning on remaining residents unless migration selection is modeled. An externally financed transfer and a tax-financed universal transfer are distinct treatments. A fiscal multiplier, if used, is a declared output change divided by the positive net fiscal injection over the same horizon; no universal multiplier follows from a meta-analysis of recipient effects. Use \(T<\infty\), \(\omega_i\ge0\), \(\sum_i\omega_i=1\), and finite means. The comparator \(p=0\) is a fully specified no-transfer fiscal regime. A household leaving the region remains in the baseline-weighted consumption target unless a distinct resident-population estimand is explicitly chosen. Financing imposes a government intertemporal budget, including debt repayment and displaced expenditures; policies with the same \((b,q)\) but different \(f\) need not have equal effects. The requested joint identified set consists of the consumption, price, and output responses over a declared feasible menu and horizons, under an admissible equilibrium and assignment law generating the observed data.
\subsection*{Variants and assumption changes}
\begin{enumerate}
\item Saturation experiment (editorial): randomize village or region coverage and recipient status jointly, identifying direct and spillover effects on prices and production at the studied scale.
\item Financing and persistence (editorial): compare equal gross transfers financed by outside aid versus domestic taxes and assess consumption after withdrawal. Existing meta-analysis estimates are background evidence; economy-wide fiscal closure is additional structure.
\end{enumerate}
\subsection*{References and source audit}
\textbf{Support and access.} Original metadata checked and 2024 primary PDF read. NBER search metadata reports an August 2026 revision, whose full text was inaccessible; no claim is made to have inspected that revision. \textbf{Locator:} IPR WP-24-21, September 30, 2024 version, abstract.
\begin{enumerate}\raggedright
\item\hypertarget{definitions:source:30007654:1}{} Tommaso Crosta; Dean Karlan; Finley Ong; Julius R\"{u}schenpöhler; Christopher Udry. 2024. \emph{Unconditional Cash Transfers: A Bayesian Meta-Analysis of Randomized Evaluations in Low and Middle Income Countries}. September 30, 2024 version, abstract.
\url{https://www.ipr.northwestern.edu/documents/working-papers/2024/wp-24-21.pdf}.
\end{enumerate}

\subsection*{Common conventions: Source mathematical conventions}

These conventions apply to each entry unless explicitly replaced.

\textbf{Models and identification.} A primitive model \(m\) specifies all probability laws, feasible choices, information, equilibrium and measurement rules used by its target. The admissible class \(\mathcal M\) is fixed independently of outcomes. If \(P_O(m)\) is its observable probability law and \(\tau(m)\) the target, its identified set at observable law \(P\) is
\[
 \mathcal I_\tau(P)=\{\tau(m):m\in\mathcal M,\ P_O(m)=P\}.
\]
Point identification means that this set is a singleton. An empty set signals incompatibility of data and assumptions. Coordinate bounds need not describe the joint set. Bounds are sharp only if every retained target is attainable or arbitrarily approachable by compatible models. Finite bounds require explicit restrictions on unobserved quantities; unrestricted confounding, hidden wealth, extrapolation or arbitrary equilibrium selection can yield uninformative sets.

\textbf{Causal interventions.} A potential outcome \(Y_i(p)\) is the outcome under a complete policy regime \(p\), including timing, financing, information and market spillovers. The target population and units are fixed before treatment. With interference, use the complete assignment vector or a justified exposure mapping; individual no-interference notation alone cannot identify an economy-wide intervention. A difference of mean potential outcomes does not require the cross-world joint distribution. Conditioning on post-treatment migration, survival, enrollment or employment changes the estimand and needs separate justification.

\textbf{Statistical statements.} An estimator, confidence set, forecast or test specifies a sampling law, model class and loss/coverage criterion. Uniform \((1-\alpha)\)-coverage means \(\inf_{m\in\mathcal M}\Pr_m\{\tau(m)\in C_n\}\geq1-\alpha\), or an explicitly asymptotic version. The whole parameter space trivially has coverage, so informativeness requires a stated width, risk or power objective. A forecasting improvement is relative to a named benchmark, information set and evaluation population.

\textbf{Equilibrium and optimization.} An equilibrium comprises feasible plans, optimality, beliefs where needed, budget constraints and all market/resource clearing. The primitives or a stated benchmark define each correspondence \(\mathcal E(p;m)\). Nonemptiness is assumed or proved, never inferred just from compact policy sets. A selected-equilibrium target uses a declared selection rule; a robust target quantifies over all permitted equilibria. Use suprema or infima unless attainment is established. An \(\varepsilon\)-optimal policy is feasible and within \(\varepsilon\) of the finite optimum value. Report an empty feasible set separately.

\textbf{Welfare and accounting.} Normative utility, interpersonal normalization, social weights, numeraire, discounting and the use of public receipts are primitives. Behavioral utility need not coincide with the welfare criterion. Household welfare includes ownership income and rents once; adding profits again to compensating variation generally double-counts them. A monetary equivalent is defined by an explicit transfer and price convention, with monotonicity and range conditions. Average effects, mechanism contributions, transfers and welfare losses are different targets. Infinite sums require convergence and dynamic feasible plans require terminal or no-Ponzi conditions.

\textbf{Source versions.} A working-paper date, revision date and journal publication year are distinguished. A DOI for a journal version is not automatically the DOI of an earlier manuscript. Locators refer to the stated version and printed pages where specified. A metadata-only check does not verify a claim about a full-text conclusion. Every entry's source subsection narrows the provenance claim accordingly.
% END ECONOMICS STATEMENT
\end{document}
